\section{Related Work}

\subsection{Model-Heterogeneous Federated Learning}

Clients with different architectures cannot be combined by direct parameter averaging. HFL methods instead exchange predictions, representations, or other architecture-independent knowledge. FedDF distills an ensemble of client predictions on unlabeled data \citep{lin2020feddf}. FedProto communicates class prototypes, and FedTGP replaces direct prototype averaging with trainable global prototypes and an adaptive-margin server objective \citep{tan2022fedproto,zhang2024fedtgp}. RHFL uses confidence-weighted collaboration to handle noisy labels \citep{fang2022rhfl}. AugHFL targets common corruptions under model heterogeneity; RAHFL extends it with diversity-enhanced local representation learning and asymmetric collaboration \citep{fang2023aughfl,fang2025rahfl}. These methods motivate the communication and robustness baselines in our study. Our question is narrower: whether a client-specific association between class and corruption induces a directional shortcut, and whether that response is formed locally or through collaboration.

\subsection{Corruption Robustness and Consistency Learning}

CIFAR-C and ImageNet-C evaluate robustness to common noise, blur, weather, and digital distortions \citep{hendrycks2019corruptions}. AugMix improves corruption robustness by combining stochastic augmentation chains with Jensen--Shannon consistency \citep{cubuk2020augmix}. Such consistency acts around views of an observed sample. CLE-HFL changes the data-generating structure by associating operators with task classes during client training. A predictor can therefore be consistent within an augmentation neighborhood yet respond differently across operators that carry different class associations. Our paired evaluation targets this cross-operator response rather than corruption robustness in general.

\subsection{Spurious Correlation, Group Robustness, and Environment Inference}

Shortcut learning describes predictive rules that work under the training distribution but fail when incidental correlations change \citep{geirhos2020shortcut}. GroupDRO minimizes worst-group risk when group identities are available \citep{sagawa2020groupdro}. GEORGE estimates hidden subclasses before applying group-robust optimization \citep{sohoni2020george}; EIIL infers environments that expose invariance violations without a predefined taxonomy \citep{creager2021eiil}; and JTT upweights examples misclassified by an initial ERM model \citep{liu2021jtt}. CVaR and related DRO formulations emphasize high-loss tails \citep{rockafellar2000cvar,duchi2021uniform}. PEW does not claim taxonomy-free environment discovery: it predicts a fixed, interpretable corruption-family taxonomy learned from public synthetic supervision. BER then changes class-conditional risk mass over those pseudo-environments. We use protocol-matched CVaR-DRO and shared-PEW GroupDRO controls to test whether generic tail emphasis or access to pseudo-groups already explains the observed behavior.

\subsection{Federated Spurious Learning and Counterfactual Diagnosis}

Spurious features are not new to federated learning. Local personalization can reintroduce client-specific shortcuts \citep{wang2024pflspurious}; FedPIN studies shortcut-averse personalized invariant learning \citep{tang2024fedpin}; and FedCD considers spurious correlation in federated domain generalization \citep{ma2024fedcd}. Counterfactual shortcut evaluation also has precedents outside FL, including medical imaging \citep{vigneshwaran2026counterfactual,jabbour2020chestxray}. CLE-HFL occupies a specific intersection of these lines: model heterogeneity, client-dependent class--corruption bindings, source-paired operator intervention, and a matched local-versus-communication attribution. Our scope is this joint problem and evidence structure; none of the individual ingredients is presented as a first use.
